\documentclass[11pt,a4paper]{article}

\usepackage[utf8]{inputenc}
\usepackage[T1]{fontenc}
\usepackage[spanish,es-tabla]{babel}
\usepackage{lmodern}
\usepackage[a4paper,margin=2.2cm]{geometry}
\usepackage[table]{xcolor}
\usepackage{booktabs}
\usepackage{enumitem}
\usepackage[most]{tcolorbox}
\usepackage{microtype}
\usepackage{amssymb}
\usepackage{titlesec}
\usepackage[hidelinks]{hyperref}

\definecolor{accent}{HTML}{1F6FEB}
\definecolor{accent2}{HTML}{0A8754}
\definecolor{warn}{HTML}{D1351A}
\definecolor{bgbox}{HTML}{EDF2F7}
\definecolor{ink}{HTML}{1F2328}
\definecolor{muted}{HTML}{57606A}

\titleformat{\section}{\Large\bfseries\color{accent}}{Bloque \thesection.}{0.6em}{}
\titleformat{\subsection}{\large\bfseries\color{ink}}{Slide \thesubsection.}{0.5em}{}
\titlespacing*{\subsection}{0pt}{1.2em}{0.4em}

\newtcolorbox{decir}{
  enhanced, breakable,
  colback=white, colframe=accent,
  boxrule=0pt, leftrule=3pt,
  arc=2pt, outer arc=0pt,
  left=8pt, right=8pt, top=4pt, bottom=4pt,
  title={\textbf{\textcolor{accent}{Qué decir}}},
  fonttitle=\small,
  attach title to upper={\par\medskip}
}

\newtcolorbox{nota}{
  enhanced, breakable,
  colback=bgbox, colframe=accent2,
  boxrule=0pt, leftrule=3pt,
  arc=2pt, outer arc=0pt,
  left=8pt, right=8pt, top=4pt, bottom=4pt,
  title={\textbf{\textcolor{accent2}{Info de respaldo (de la nota)}}},
  fonttitle=\small,
  attach title to upper={\par\medskip}
}

\newtcolorbox{transicion}{
  enhanced, breakable,
  colback=bgbox, colframe=muted,
  boxrule=0pt, leftrule=3pt,
  arc=2pt, outer arc=0pt,
  left=8pt, right=8pt, top=4pt, bottom=4pt,
  title={\textbf{\textcolor{muted}{Transición}}},
  fonttitle=\small,
  attach title to upper={\par\medskip}
}

\newcommand{\nref}[1]{\textcolor{muted}{\small(\texttt{#1})}}

\title{\textbf{Guión de presentación — TP3 Perceptrones}\\
       \large 35 slides · qué decir en cada una}
\author{}
\date{}

\begin{document}
\maketitle

\noindent\textcolor{muted}{\small Cada entrada corresponde al número de slide del PDF
\texttt{slides\_presentacion.pdf}. Las cajas verdes incluyen información de respaldo
extraída de las notas en \texttt{docs/notas/} para responder preguntas. La presentación
arranca con Ej 1 (no hay sección de ``Decisiones de diseño'' al inicio: las decisiones
se discuten dentro del bloque de Ej 2 una vez que ya está planteado el problema).}

\tableofcontents
\vspace{0.5em}
\noindent\textcolor{muted}{\small Tiempo estimado total: 18--22 minutos.}

% =====================================================================
\section{Apertura}

\subsection{Slide 1 — Portada}

\begin{decir}
\begin{itemize}[leftmargin=*,topsep=0pt]
  \item Buenas, somos \emph{[nombres del grupo]}.
  \item Vamos a presentar nuestros resultados experimentales sobre tres ejercicios:
        detección de fraude por knowledge distillation, clasificación de dígitos
        manuscritos, y mejora de esa clasificación.
  \item La presentación está organizada en tres bloques, uno por ejercicio.
\end{itemize}
\end{decir}

% =====================================================================
\section{Ejercicio 1 — Knowledge Distillation}

\subsection{Slide 2 — Divider Ej 1}

\begin{decir}
\begin{itemize}[leftmargin=*,topsep=0pt]
  \item Empezamos con el Ej 1: \textbf{detección de fraude por knowledge distillation}.
  \item Es un perceptrón simple, pensado como prueba de concepto.
\end{itemize}
\end{decir}

\subsection{Slide 3 — Ej 1: el problema}

\begin{decir}
\begin{itemize}[leftmargin=*,topsep=0pt]
  \item \emph{Knowledge distillation} = replicar la salida de un BigModel pre-entrenado
        con un TinyModel (acá, un perceptrón simple).
  \item El target de entrenamiento es \texttt{big\_model\_fraud\_probability}, una
        \textbf{probabilidad continua entre 0 y 1}.
  \item \textbf{Regla clave del enunciado}: \texttt{flagged\_fraud} (la etiqueta binaria
        real) \textbf{nunca se usa para entrenar}; sólo aparece después como criterio
        de evaluación.
  \item Como la salida tiene que estar en $[0,1]$, la activación natural es la
        \textbf{sigmoide}.
\end{itemize}
\end{decir}

\subsection{Slide 4 — Ej 1: análisis del dataset (tres reglas determinísticas)}

\begin{decir}
\begin{itemize}[leftmargin=*,topsep=0pt]
  \item \textbf{Antes de entrenar nada} hicimos exploración del dataset.
  \item Encontramos tres features con \textbf{umbrales duros} que separan
        perfectamente las clases (todo lo de la derecha del umbral es fraude):
        \texttt{amount\_usd > 500}, \texttt{quantity $\geq$ 10}, \texttt{items\_viewed $\geq$ 15}.
  \item Cada regla individual marca cientos de filas con \textbf{cero falsos positivos}.
  \item Aplicadas como OR, dan \textbf{TP=695, FP=0, FN=174} sobre 7.500 filas:
        \emph{precision 100\%, recall 80\%, accuracy 97.7\%}.
\end{itemize}
\end{decir}

\begin{nota}
\textbf{Por qué importa el análisis previo del dataset:} Pearson y Spearman cercanos
a 0 en \texttt{timestamp}, \texttt{device\_screen\_resolution} y
\texttt{time\_since\_last\_login\_s} indican que esas tres features no aportan
señal lineal ni monótona. Las descartamos para el perceptrón simple del Ej 1
porque, además, dos de ellas tienen magnitudes de orden $10^9$ y $10^6$ que
dominarían el cálculo del gradiente sin normalización perfecta.
\nref{decisiones\_y\_backing\_teorico.md}
\end{nota}

\subsection{Slide 5 — Ej 1: lineal vs no-lineal (MSE)}

\begin{decir}
\begin{itemize}[leftmargin=*,topsep=0pt]
  \item Mismo dataset, K-fold = 5 estratificado, mismas features. La única diferencia
        entre los dos modelos es la \textbf{presencia de la sigmoide}.
  \item El no-lineal saca \textbf{59\% menos MSE}: $0.0110 \pm 0.0004$ vs
        $0.0266 \pm 0.0008$.
  \item \textbf{¿Por qué gana tan claro?} El target es una probabilidad acotada en
        $[0,1]$. La sigmoide \emph{naturalmente} mapea su salida a ese rango. El
        lineal puede producir valores fuera de $[0,1]$ y el MSE penaliza fuerte
        esa salida fuera de dominio.
\end{itemize}
\end{decir}

\subsection{Slide 6 — Ej 1: pruebas sobre threshold}

\begin{decir}
\begin{itemize}[leftmargin=*,topsep=0pt]
  \item Con el threshold default $0.5$ las métricas se ven \emph{horribles}
        (precision $0.33$, F1 $0.50$). \textbf{No es el modelo, es el threshold}.
  \item Las salidas de la sigmoide se concentran arriba: la mediana de los scores
        de fraude es $0.99$. Los no-fraudes \emph{no} bajan a $0.05$, bajan a
        $0.4$--$0.6$. Por eso el $0.5$ sobre-detecta.
  \item Hicimos un sweep y el \textbf{óptimo por F1 es $0.89$}: precision $0.886$,
        recall $0.861$, F1 $0.873$, accuracy $0.971$.
  \item El modelo aprendió un buen \textbf{ranking}, no una calibración perfecta.
\end{itemize}
\end{decir}

\subsection{Slide 7 — Ej 1: recomendación de umbral}

\begin{decir}
\begin{itemize}[leftmargin=*,topsep=0pt]
  \item La elección de threshold depende del \textbf{costo asimétrico} entre falso
        positivo y falso negativo.
  \item Comparamos cuatro thresholds en la tabla: el default $0.5$ deja 1743 FP;
        el óptimo F1 ($0.89$) deja 96 FP; el $0.99$ deja \textbf{cero} FP a costa
        de bajar el recall a $0.527$.
  \item Conclusión: \textbf{el modelo entrega scores; la política de threshold es
        decisión de negocio}.
\end{itemize}
\end{decir}

% =====================================================================
\section{Ejercicio 2 — Clasificación de dígitos}

\subsection{Slide 8 — Divider Ej 2}

\begin{decir}
\begin{itemize}[leftmargin=*,topsep=0pt]
  \item Pasamos al Ej 2: \textbf{clasificación multiclase} de dígitos manuscritos
        con un MLP.
  \item Acá entran las decisiones de diseño que vamos a discutir en bloque.
\end{itemize}
\end{decir}

\subsection{Slide 9 — Ej 2: el problema}

\begin{decir}
\begin{itemize}[leftmargin=*,topsep=0pt]
  \item Diez clases (dígitos 0-9). Train: 12.449 muestras. Test: 2.498.
  \item Cada muestra: vector de 784 floats $\in [0,1]$ — son imágenes 28$\times$28
        \emph{flatten}.
  \item \textbf{Arquitectura final}: $[784, 100, 50, 10]$ con ReLU en las capas
        ocultas y softmax en la salida.
  \item \textbf{Regla del TP}: el test set se evalúa \emph{una sola vez al final}.
        Toda la búsqueda de hiperparámetros es sobre \texttt{digits.csv}.
\end{itemize}
\end{decir}

\subsection{Slide 10 — Ej 2: fases de experimentación}

\begin{decir}
\begin{itemize}[leftmargin=*,topsep=0pt]
  \item Workflow disciplinado en \textbf{dos fases}.
  \item \textbf{Fase 1} — \emph{exploratoria, k\_folds = 1}: sweeps secuenciales
        (arquitectura $\to$ optimizador $\to$ LR $\to$ batch size) para reducir
        el espacio de búsqueda barato.
  \item Se consolida el ganador en un \texttt{base.json}.
  \item \textbf{Fase 2} — \emph{validación, k\_folds = 5}: comparativas one-at-a-time
        para confirmar que las elecciones son robustas, no suerte de un solo split.
  \item Sólo después corre el \texttt{final\_eval} sobre \texttt{digits\_test.csv},
        \emph{una sola vez}.
\end{itemize}
\end{decir}

\subsection{Slide 11 — Ej 2: sweep de arquitectura}

\begin{decir}
\begin{itemize}[leftmargin=*,topsep=0pt]
  \item Comparamos cinco arquitecturas: la base, una con misma forma pero un poco
        más ancha ($[784, 128, 64, 10]$), una \emph{wider}, una \emph{deeper} y
        una \emph{shallow}.
  \item \textbf{Empate técnico} entre la base y $[784, 128, 64, 10]$
        ($\Delta < 0.4\%$, std $\sim 0.4\%$).
  \item \textbf{Wider} ($[784, 200, 100, 10]$) mejora val (+0.17 pp) pero
        \emph{empeora test} ($-0.48$ pp) y es $3.6\times$ más lento.
  \item \textbf{Decisión por parsimonia (regla de Occam)}: elegimos la más simple
        y rápida. \emph{Wider} es el ejemplo empírico de la curva U: aumentar
        capacidad sin más datos $\Rightarrow$ overfitting.
\end{itemize}
\end{decir}

\begin{nota}
\textbf{Capacidad del modelo (slide 8 del PDF de regularización):} ``habilidad
de un modelo de aproximar una variedad de funciones''. Cambia con: elección de
modelo, arquitectura, cantidad de features de input. La curva U del slide 9
muestra que hay un punto óptimo: por debajo, underfitting; por encima,
overfitting. \emph{Wider} sin más datos cae en la zona de overfitting — es lo
que vimos. \nref{decisiones\_y\_backing\_teorico.md \S 6}
\end{nota}

\subsection{Slide 12 — Ej 2: sweep de optimizador}

\begin{decir}
\begin{itemize}[leftmargin=*,topsep=0pt]
  \item Tres optimizadores: SGD vanilla, Momentum ($\beta=0.9$), Adam.
  \item \textbf{Adam gana}: $96.74\%$, best\_epoch $= 7$.
  \item Momentum queda dentro del desvío. SGD vanilla \emph{no convergió en 50 epochs}.
  \item Elegimos Adam por convergencia más rápida y resultado igual o mejor.
\end{itemize}
\end{decir}

\subsection{Slide 13 — Función de activación}

\begin{decir}
\begin{itemize}[leftmargin=*,topsep=0pt]
  \item Por contexto: sigmoide en Ej 1 (target $[0,1]$); \textbf{ReLU} en capas
        ocultas de Ej 2/3; \textbf{softmax} en la salida (clasificación multiclase).
  \item \textbf{¿Qué es saturar?} Una activación satura cuando $f'(z) \to 0$
        para entradas grandes. Sigmoide y tanh saturan en sus extremos: el
        gradiente que llega a las capas previas se desvanece (\emph{vanishing
        gradient}) y la red deja de aprender.
  \item \textbf{ReLU $= \max(0, z)$}: derivada $1$ en zona positiva, $0$ en
        negativa. Sin saturación en su lado activo $\Rightarrow$ gradiente
        constante a través de muchas capas.
\end{itemize}
\end{decir}

\begin{nota}
\textbf{ReLU en detalle:}
\begin{itemize}[leftmargin=*,topsep=2pt]
  \item Con buena inicialización, $\sim$50\% de las neuronas están apagadas en
        cada batch. Es \emph{feature, no bug} (sparsity).
  \item \emph{Neuronas muertas}: si una neurona devuelve 0 para todos los
        datos, su gradiente queda en 0 y nunca se actualiza. Suele pasar con
        learning rate alto o mala init.
  \item \textbf{He init} ($\sigma = \sqrt{2/\text{fan\_in}}$) compensa que ReLU
        ``mata'' la mitad del input — mantiene varianza estable a través de
        capas. Por eso usamos He cuando hay ReLU y Xavier cuando hay
        tanh/sigmoid.
\end{itemize}
\nref{relu.md}
\end{nota}

\subsection{Slide 14 — Inicialización de pesos}

\begin{decir}
\begin{itemize}[leftmargin=*,topsep=0pt]
  \item \textbf{¿Por qué no cero?} Si todas las neuronas arrancan iguales, hacen
        lo mismo en cada capa. El gradiente las mueve igual — \emph{nunca se
        diferencian}. Es el problema de simetría.
  \item \textbf{¿Por qué pequeños?} Pesos grandes saturan las activaciones.
  \item Tres esquemas: \textbf{He} ($\sigma = \sqrt{2/\text{fan\_in}}$) para
        ReLU; \textbf{Xavier} ($\sigma = \sqrt{1/\text{fan\_in}}$) para
        tanh/sigmoid/softmax; \textbf{Uniforme} $[-0.1, 0.1]$ para el perceptrón
        simple del Ej 1.
\end{itemize}
\end{decir}

\begin{nota}
\textbf{De dónde vienen las fórmulas:} la cátedra mandó un HTML sobre weight
initialization que cubre Xavier (Glorot \& Bengio 2010) y He (He et al. 2015).
Nuestra fórmula de Xavier usa $\sqrt{1/\text{fan\_in}}$ (variante LeCun); el
HTML cita la forma Glorot $\sqrt{2/(\text{n\_in}+\text{n\_out})}$. \emph{Si
preguntan}: ambas formas existen en la literatura (PyTorch tiene las dos),
empíricamente quedaron en empate técnico ($\sim$96\%, std $\sim 0.3\%$).
\nref{decisiones\_y\_backing\_teorico.md \S 2}
\end{nota}

\subsection{Slide 15 — Optimizador}

\begin{decir}
\begin{itemize}[leftmargin=*,topsep=0pt]
  \item \textbf{SGD vanilla}: $w \leftarrow w - \eta \nabla L$. Lo más simple. Lento
        en valles alargados.
  \item \textbf{Momentum} ($\beta = 0.9$): acumula \emph{velocidad} en la dirección
        del gradiente; atraviesa valles sin oscilar.
  \item \textbf{Adam}: combina momentum con \emph{learning rate adaptativo por
        parámetro}.
  \item Resultado en Ej 2 lo vimos en el sweep: Adam $\approx$ Momentum $\gg$ SGD.
        Elegimos Adam por convergencia más rápida.
\end{itemize}
\end{decir}

\begin{nota}
\textbf{Adam = Momentum + RMSProp:} mantiene un promedio móvil del gradiente
(como momentum) y otro del gradiente al cuadrado (como RMSProp). Divide el
update por la raíz del segundo momento $\Rightarrow$ learning rate adaptativo
por parámetro. Defaults: $\alpha = 10^{-3}$, $\beta_1 = 0.9$, $\beta_2 = 0.999$,
$\varepsilon = 10^{-8}$. \nref{decisiones\_y\_backing\_teorico.md \S 3}
\end{nota}

\subsection{Slide 16 — Learning rate}

\begin{decir}
\begin{itemize}[leftmargin=*,topsep=0pt]
  \item Recomendación de la clase: learning rate \emph{pequeño} y probar varios
        órdenes de magnitud.
  \item Demasiado grande $\to$ la red oscila o diverge. Demasiado chico $\to$
        no converge en tiempo razonable.
  \item Hicimos sweep en \textbf{dos pasos}: Fase 1 con cinco valores razonables;
        Fase 2 con extremos catastróficos para mapear los bordes.
  \item Las próximas dos slides muestran las tablas de cada fase.
\end{itemize}
\end{decir}

\subsection{Slide 17 — Learning rate: Fase 1 exploratoria (k=1)}

\begin{decir}
\begin{itemize}[leftmargin=*,topsep=0pt]
  \item \textbf{$lr = 10^{-3}$} ganó claro: $96.74\%$, best\_epoch $= 7$.
  \item $lr = 10^{-4}$ converge pero \emph{muy lento} (best\_ep $= 47$).
  \item $lr \geq 5 \cdot 10^{-3}$ corta a las 2-4 épocas: \textbf{overshoot}.
  \item Tres criterios apuntan al mismo número: máximo de val\_acc, convergencia
        rápida y ausencia de overshoot.
\end{itemize}
\end{decir}

\begin{nota}
\textbf{Tabla de Fase 1.3 completa} (5 valores, k=1):
\begin{itemize}[leftmargin=*,topsep=2pt]
  \item $0.001 \to 96.74\%$, best\_ep=7. $\checkmark$ Óptimo.
  \item $0.0005 \to 96.54\%$, best\_ep=17. Lento.
  \item $0.0001 \to 96.22\%$, best\_ep=47. No converge.
  \item $0.005 \to 96.02\%$, best\_ep=4. Overshoot moderado.
  \item $0.01 \to 95.30\%$, best\_ep=2. Overshoot fuerte.
\end{itemize}
\nref{lr\_sweep\_conclusion.md}
\end{nota}

\subsection{Slide 18 — Learning rate: Fase 2 validación (k=5)}

\begin{decir}
\begin{itemize}[leftmargin=*,topsep=0pt]
  \item Mismo ganador, ahora con K-fold = 5: \textbf{$10^{-3}$ confirmado
        ($96.22\%$)}.
  \item Los \textbf{extremos catastróficos} confirman la teoría:
        $lr = 0.1$ diverge ($29\%$), $lr = 10$ explota a NaN — la red queda
        en accuracy random ($\sim 12\%$, equivalente a $1/10$ clases).
  \item Robusto: dos sweeps independientes coinciden.
\end{itemize}
\end{decir}

\subsection{Slide 19 — Convergencia y bias}

\begin{decir}
\begin{itemize}[leftmargin=*,topsep=0pt]
  \item Dos criterios distintos por la naturaleza de cada problema. \textbf{Ej 1}:
        $\varepsilon$ absoluto sobre MSE de train (regla del apunte para
        perceptrón simple). \textbf{Ej 2/3}: \emph{early stopping} sobre
        \texttt{val\_loss} con \texttt{patience = 10}.
  \item Early stopping detecta overfitting que $\varepsilon$ solo no detecta.
  \item \textbf{Bias trick}: agregamos columna de 1's al input de cada capa,
        los pesos quedan con shape \texttt{(n\_out, n\_in+1)}. El bias es un
        peso más, se actualiza con la misma regla.
\end{itemize}
\end{decir}

\begin{nota}
\textbf{Early stopping (slide 14 del PDF de regularización):}
\begin{itemize}[leftmargin=*,topsep=2pt]
  \item En cada época trackeamos \emph{best\_val\_loss} y guardamos los pesos del
        mejor modelo.
  \item Si la val\_loss no mejora durante \texttt{patience} épocas seguidas,
        cortamos y \emph{restauramos los pesos del mejor modelo} (no del último).
  \item \texttt{epochs} máximo en Ej 2/3 = 50 con patience = 10. \emph{best\_epoch}
        promedio $\sim 5$ — la red converge mucho antes de los 50.
\end{itemize}
\nref{early\_stopping.md}
\end{nota}

\subsection{Slide 20 — Ej 2: curvas de aprendizaje (K=5)}

\begin{decir}
\begin{itemize}[leftmargin=*,topsep=0pt]
  \item \textbf{Convergencia limpia} en $\sim 5$ epochs.
  \item Train y validación bajan \emph{juntos}: \textbf{sin overfitting visible}.
  \item Las bandas son los desvíos del K-fold; estrechas $\Rightarrow$ entrenamiento
        estable entre folds.
\end{itemize}
\end{decir}

\subsection{Slide 21 — Ej 2: matriz de confusión del base (val K=5)}

\begin{decir}
\begin{itemize}[leftmargin=*,topsep=0pt]
  \item Diagonal limpia para \textbf{9 de 10} clases.
  \item \textbf{La clase 8 colapsa}: precision = recall = 0.
  \item Hint fuerte: algo raro pasa con la \textbf{distribución de los datos}
        (cobertura insuficiente del 8 en \texttt{digits.csv}).
\end{itemize}
\end{decir}

\subsection{Slide 22 — Ej 2: final eval, el cachetazo}

\begin{decir}
\begin{itemize}[leftmargin=*,topsep=0pt]
  \item Acá viene la sorpresa.
  \item val K-fold $= 96.22\%$. \textbf{Test} (\texttt{digits\_test.csv})
        $= 86.30\%$. \textbf{Drop de 10 puntos porcentuales}.
  \item \textbf{No es underfitting, no es overfitting clásico}: train y val bajaron
        juntos. Es \textbf{distribution shift} — los datos de test son una
        población distinta de los de train.
\end{itemize}
\end{decir}

\begin{nota}
\textbf{¿Qué es distribution shift?} La distribución de los datos de
entrenamiento es distinta de la de evaluación: el modelo aprende patrones
de train que no generalizan al test, no por memorizar (overfitting) sino
porque \emph{train y test son poblaciones distintas}.
\begin{itemize}[leftmargin=*,topsep=2pt]
  \item \emph{Covariate shift}: $P_{\text{train}}(x) \neq P_{\text{test}}(x)$,
        pero $P(y|x)$ es la misma. Es lo que vimos acá.
  \item Síntomas en el TP3: drop val$\to$test asimétrico, clase 8 colapsada,
        errores distribuidos (no concentrados).
\end{itemize}
\nref{explicacion\_distribution\_shift.md}
\end{nota}

\subsection{Slide 23 — Ej 2: matriz sobre digits\_test.csv}

\begin{decir}
\begin{itemize}[leftmargin=*,topsep=0pt]
  \item Las confusiones están \textbf{distribuidas}, no concentradas en una clase.
        Ese patrón es síntoma de \emph{shift}, no de falta de capacidad
        (capacidad insuficiente daría errores concentrados).
  \item Conclusión: \textbf{más datos, no más modelo}. Eso es lo que el Ej 3
        provee con \texttt{more\_digits.csv}.
\end{itemize}
\end{decir}

% =====================================================================
\section{Ejercicio 3 — Mejorar la clasificación}

\subsection{Slide 24 — Divider Ej 3}

\begin{decir}
\begin{itemize}[leftmargin=*,topsep=0pt]
  \item Cerramos con el Ej 3: \textbf{llevar el accuracy de test al 98\%}.
  \item Vamos a probar si la hipótesis del shift es correcta.
\end{itemize}
\end{decir}

\subsection{Slide 25 — Ej 3: hipótesis y plan}

\begin{decir}
\begin{itemize}[leftmargin=*,topsep=0pt]
  \item \textbf{Hipótesis}: el techo del Ej 2 es por distribution shift, no por
        falta de capacidad.
  \item Plan secuencial: (1) \textbf{más datos} con \texttt{more\_digits.csv}
        (15.742 muestras adicionales) — \emph{cero cambios al modelo}; (2) si
        no llega al 98\%, agregamos regularización (L2, dropout, augmentation).
\end{itemize}
\end{decir}

\subsection{Slide 26 — Ej 3: el movimiento dominó}

\begin{decir}
\begin{itemize}[leftmargin=*,topsep=0pt]
  \item Con \emph{sólo} concatenar \texttt{more\_digits.csv}:
        \textbf{val $96.22 \to 97.06$}, \textbf{test $86.30 \to 96.36$}.
  \item \textbf{+10 puntos porcentuales en test}, \emph{cero cambios al modelo}.
  \item Macro F1 $0.857 \to 0.959$. La clase 8 que estaba colapsada en Ej 2
        ahora se predice correctamente.
  \item \textbf{Confirma la hipótesis del shift}: era cobertura, no capacidad.
\end{itemize}
\end{decir}

\subsection{Slide 27 — Ej 3: matriz de confusión del ganador (test)}

\begin{decir}
\begin{itemize}[leftmargin=*,topsep=0pt]
  \item Diagonal mucho más limpia.
  \item Los errores residuales se reparten — síntoma de que \textbf{ya no hay un
        agujero específico de cobertura}, sino limitaciones globales del modelo.
\end{itemize}
\end{decir}

\subsection{Slide 28 — Ej 3: resultados}

\begin{decir}
\begin{itemize}[leftmargin=*,topsep=0pt]
  \item Probamos siete configuraciones tomando \texttt{base\_extra\_data} como
        referencia (test $96.36$).
  \item \textbf{L2} ($\lambda = 10^{-4}$): $+0.20$ pp consistente.
  \item \textbf{Dropout} ($p = 0.2$): $-0.08$ pp en test (gana en val, pierde en test).
  \item \textbf{Ganador: L2 + augmentation gaussiana ($\sigma = 0.05$) $\to$ 96.88\%}
        en test ($+0.52$ pp sobre el baseline con extra data).
  \item \emph{Wider + L2 + aug} no supera al ganador: descarta falta de capacidad.
\end{itemize}
\end{decir}

\begin{nota}
\textbf{Resumen de cada técnica} (qué hace + cómo se implementa):
\begin{itemize}[leftmargin=*,topsep=2pt]
  \item \textbf{L2}: agrega $\frac{1}{2} \lambda \|w\|^2$ al loss; gradiente
        $+\lambda w$. \emph{El bias no se penaliza} (la columna 0 de cada matriz).
  \item \textbf{Augmentation gaussiana}: en cada batch suma ruido
        $\mathcal{N}(0, \sigma^2)$ a las entradas. $\sigma = 0.05$ es lo que
        ganó.
  \item \textbf{Dropout}: cada neurona se ``apaga'' con probabilidad $p$ durante
        training; en inferencia se usa la red entera. Inverted dropout: la
        división por $(1-p)$ compensa la baja densidad de activaciones.
\end{itemize}
\nref{resultados\_ejercicio3\_explicacion.md, clase\_regularizacion\_cosas\_implementadas.md}
\end{nota}

\subsection{Slide 29 — Ej 3: comparación visual con Ej 2}

\begin{decir}
\begin{itemize}[leftmargin=*,topsep=0pt]
  \item Línea roja: target $\geq 98\%$.
  \item \textbf{El salto principal viene del baseline con extra data}, no de la
        regularización fina.
  \item La regularización aporta un techo modesto ($\sim 0.5$ pp) sobre ese
        baseline.
\end{itemize}
\end{decir}

\subsection{Slide 30 — Ej 3: curvas del modelo ganador}

\begin{decir}
\begin{itemize}[leftmargin=*,topsep=0pt]
  \item Convergencia más lenta que en Ej 2 base (best\_epoch $\sim 10$ vs $5$)
        por L2 + augmentation, que ``cuestan'' alguna época de optimización a
        cambio de mejor generalización.
  \item Sin overfitting visible.
\end{itemize}
\end{decir}

\subsection{Slide 31 — Ej 3: qué ayudó y qué no}

\begin{decir}
\begin{itemize}[leftmargin=*,topsep=0pt]
  \item \textbf{Ayudó}: más datos ($+10$ pp, dominante), L2 ($+0.20$ pp),
        aug gaussiana ($+0.32$ pp adicional sobre L2 — \emph{sinergia}).
  \item \textbf{No ayudó}: dropout (gana en val, pierde en test — reduce
        varianza al fold, no al shift), wider (mejor val, peor test —
        confirma que no falta capacidad), combinar todo (L2+dropout+aug peor
        que L2+aug solo).
\end{itemize}
\end{decir}

\begin{nota}
\textbf{Por qué dropout no transfirió:} reduce la varianza al fold de
cross-validation (mejora val K-fold) pero \emph{no corrige distribution shift}.
Si train y test son poblaciones distintas, dropout no pone ningún caso de test
en la mesa de entrenamiento. Lección: una técnica que mejora val no
necesariamente mejora test cuando hay shift.
\nref{explicacion\_distribution\_shift.md}
\end{nota}

\subsection{Slide 32 — Ej 3: por qué no llegamos al 98\%}

\begin{decir}
\begin{itemize}[leftmargin=*,topsep=0pt]
  \item Mejor test: $96.88\%$. Faltan $1.12$ pp.
  \item \textbf{Hipótesis principal}: el augmentation \textbf{gaussiano} simula
        ruido isotrópico por píxel. Pero el shift parece ser \textbf{geométrico}
        (rotación, traslación, grosor de trazo distintos por estilo de escritura).
  \item En la clase de regularización vimos \textbf{cuatro} formas de
        augmentation (gaussiano, rotaciones, traslaciones, escalas); sólo
        implementamos la primera.
  \item \emph{Wider} empeoró $\Rightarrow$ no es problema de capacidad. Es
        problema de cobertura.
  \item \textbf{Próximo paso natural}: augmentation geométrico.
\end{itemize}
\end{decir}

\begin{nota}
\textbf{Matching con el PDF de regularización:}
\begin{itemize}[leftmargin=*,topsep=2pt]
  \item Slide 14 — Early stopping: \emph{implementado} y activo en Ej 2/3.
  \item Slide 18 — Augmentation: gaussiano (\emph{implementado, $\sigma=0.05$}),
        rotaciones / traslaciones / escalas (\emph{NO implementados}).
  \item Slides 20-25 — L2 / Weight Decay: \emph{implementado, $\lambda = 10^{-4}$}.
  \item Slide 26 — Dropout: \emph{implementado, descartado por no transferir a test}.
\end{itemize}
Las tres formas de augmentation que NO implementamos son la hipótesis
principal del techo en 96.88\%.
\nref{clase\_regularizacion\_cosas\_implementadas.md}
\end{nota}

% =====================================================================
\section{Cierre}

\subsection{Slide 33 — Gracias}

\begin{decir}
\begin{itemize}[leftmargin=*,topsep=0pt]
  \item Gracias. Quedamos atentos a las preguntas.
\end{itemize}
\end{decir}

\end{document}
